\documentclass[11pt]{article}

\usepackage[margin=1in]{geometry}
\usepackage{amsmath}
\usepackage{amssymb}
\usepackage{amsthm}

\newtheorem{theorem}{Theorem}
\newtheorem{lemma}[theorem]{Lemma}
\newtheorem{proposition}[theorem]{Proposition}
\newtheorem{corollary}[theorem]{Corollary}
\theoremstyle{definition}
\newtheorem{conjecture}[theorem]{Conjecture}
\newtheorem{remark}[theorem]{Remark}

\newcommand{\xor}{\oplus}
\newcommand{\abs}[1]{\lvert #1 \rvert}
\setlength{\emergencystretch}{3em}

\title{Refutation of Conjecture 00000001283:\\
the nim-equation $x \xor y = x \cdot y$ has no solutions with $x, y \ge 1$}
\author{earthking11}
\date{13 September 2026}

\begin{document}
\maketitle

\begin{abstract}
We disprove Conjecture 00000001283.  In the stated domain $x, y \ge 1$ the
nim-equation $x \xor y = x \cdot y$ (with $\xor$ the carry-free binary
addition, i.e.\ bitwise XOR) has \emph{no} solutions at all.  Moreover the
pair $(2,2)$ that the conjecture lists as a solution is not one, because
$2 \xor 2 = 0 \ne 4 = 2 \cdot 2$, and the second clause of the conjecture --
that for all other pairs $x, y > 2$ the difference $x \xor y - x \cdot y$ is a
sign flip of a power of two -- fails already at $(3,3)$, where the difference
is $-9$.  Over the box $x, y \in [3, 499]$ exactly $39$ of the $247\,009$
pairs satisfy that clause; the remaining $246\,970$ violate it.  All three
facts are verified by exhaustive computation and, independently, in Lean~4.
\end{abstract}

\section{The conjecture under test}\label{sec:conjecture}

From \texttt{conjectures/00000001283.md}:

\begin{quote}
\textbf{Conjecture 00000001283.} \emph{A nim-equation is an integer solution
of $x \xor y = x \cdot y$ with $x, y \ge 1$, where $\xor$ is binary addition
without carry.  Conjecture: the only solutions are $(2,2)$ and $(0,0)$; for
all other candidate pairs with $x, y > 2$, the difference between the
nim-sum and the product is a sign flip of a power of two.}
\end{quote}

Here ``binary addition without carry'' is bitwise XOR: writing $x$ and $y$ in
base~$2$, one adds each bit column independently and discards carries.  The
claim has two clauses:

\begin{enumerate}
  \item[(C1)] the solution set of $x \xor y = x \cdot y$ over non-negative
    integers is exactly $\{(2,2), (0,0)\}$;
  \item[(C2)] for every other pair with $x, y > 2$ one has
    $x \xor y - x \cdot y = \pm 2^{j}$ for some $j \ge 0$.
\end{enumerate}

We show that the statement is false as written, and in fact that it is false
in a stronger sense than the wording suggests.

\section{Part 1: $(2,2)$ is not a solution}\label{sec:twotwo}

In binary, $2 = 10_{2}$, so
\[
  2 \xor 2 = 10_{2} \xor 10_{2} = 00_{2} = 0 ,
\]
whereas the ordinary product is $2 \cdot 2 = 4$.  Hence
\[
  2 \xor 2 = 0 \ne 4 = 2 \cdot 2 ,
\]
so $(2,2)$ does \emph{not} satisfy the nim-equation.  Clause (C1) explicitly
lists $(2,2)$ as a solution, so (C1) is literally false.  The other listed
pair, $(0,0)$, is a genuine solution ($0 \xor 0 = 0 = 0 \cdot 0$) but lies
outside the stated domain $x, y \ge 1$; every theorem below concerns that
domain.

\begin{remark}
A charitable reading of ``the only solutions are $(2,2)$ and $(0,0)$'' is
``there are no solutions \emph{besides} these two,'' which would not, by
itself, assert that $(2,2)$ is a solution.  But the sentence as written says
that the only solutions \emph{are} $(2,2)$ and $(0,0)$, i.e.\ it asserts
membership; that assertion is false.  In any case Section~\ref{sec:none}
shows there are no solutions with $x, y \ge 1$, so even the charitable
reading leaves the conjecture with an empty (and hence vacuous) true half and
a false second half.
\end{remark}

\section{Part 2: there are no solutions with $x, y \ge 1$}\label{sec:none}

We prove the stronger statement announced in the abstract.

\begin{theorem}\label{thm:none}
There is no pair of integers $(x, y)$ with $x, y \ge 1$ such that
$x \xor y = x \cdot y$.
\end{theorem}

\begin{proof}
Let $x, y \ge 1$ and suppose $x \xor y = x \cdot y$.

\emph{Case $x = 1$ (or, symmetrically, $y = 1$).}  Adding $1$ to an integer
without carry changes only the lowest bit: it toggles that bit.  Hence for
$y \ge 1$,
\[
  1 \xor y =
  \begin{cases}
    y + 1, & y \text{ even},\\
    y - 1, & y \text{ odd},
  \end{cases}
\]
so $1 \xor y = y \pm 1$, which is never equal to $1 \cdot y = y$ because
$\pm 1 \ne 0$.  Thus $x = 1$ gives no solution, and by symmetry neither does
$y = 1$.

\emph{Case $x, y \ge 2$.}  For the XOR we use the standard identity
\[
  x + y = (x \xor y) + 2\,(x \mathbin{\&} y),
\]
which holds because addition decomposes into the carry-free sum and twice the
carry bits.  Since $x \mathbin{\&} y \ge 0$,
\begin{equation}\label{eq:xorle}
  x \xor y = x + y - 2\,(x \mathbin{\&} y) \le x + y .
\end{equation}
On the other hand, for $x, y \ge 2$,
\[
  x \cdot y - (x + y) = (x - 1)(y - 1) - 1 \ge 1 \cdot 1 - 1 = 0 ,
\]
with equality if and only if $(x - 1)(y - 1) = 1$, i.e.\ $x = y = 2$.  Hence
\begin{equation}\label{eq:prodge}
  x \cdot y \ge x + y ,
\end{equation}
and combining \eqref{eq:xorle} and \eqref{eq:prodge},
\[
  x \xor y \le x + y \le x \cdot y .
\]
If $x \xor y = x \cdot y$ then both inequalities are equalities.  Equality in
\eqref{eq:prodge} forces $x = y = 2$, but then equality in \eqref{eq:xorle}
would require $x \xor y = x + y$, whereas $2 \xor 2 = 0 \ne 4 = 2 + 2$.  So
the case $x, y \ge 2$ yields no solution either.

Both cases being impossible, there is no solution with $x, y \ge 1$.
\end{proof}

\begin{corollary}
Over the non-negative integers the solution set of $x \xor y = x \cdot y$ is
exactly $\{(0,0)\}$.
\end{corollary}

\begin{proof}
$0 \xor 0 = 0 = 0 \cdot 0$ is a solution.  Conversely, if $x, y \ge 0$ solve
the equation and $(x, y) \ne (0,0)$, then either $x = 0$ or $y = 0$ (the
``boundary'' cases $x, y \ge 1$ being impossible by
Theorem~\ref{thm:none}).  If $x = 0$ then $x \xor y = y$ but
$x \cdot y = 0$, so $y = 0$; symmetrically for $y = 0$.  Hence the only
non-negative solution is $(0,0)$.
\end{proof}

This is confirmed by exhaustive search: over $0 \le x, y < 500$ the solution
set is $\{(0,0)\}$, and over $1 \le x, y < 500$ it is empty
(\texttt{reproduce.py}, checks~[2] and~[2b]; the same bounded search is
formalised in \texttt{lean4/Main.lean} for the box $0 \le x, y < 200$).

\section{Part 3: the second clause fails at $(3,3)$}\label{sec:clause2}

For a pair $(x, y)$ write
\[
  \Delta(x, y) = (x \xor y) - x \cdot y .
\]
Clause (C2) asserts that $\Delta(x, y) \in \{\, \pm 2^{j} : j \ge 0 \,\}$ for
every pair with $x, y > 2$ and $(x, y) \notin \{(2,2), (0,0)\}$.  This is
false.  At $(3,3)$,
\[
  3 \xor 3 = 11_{2} \xor 11_{2} = 00_{2} = 0,
  \qquad 3 \cdot 3 = 9,
  \qquad \Delta(3,3) = 0 - 9 = -9 .
\]
The number $9$ is not a power of two.  Indeed the powers of two are
$2^{0} = 1$, $2^{1} = 2$, $2^{2} = 4$, $2^{3} = 8$, $2^{4} = 16$, and
$2^{j}$ is strictly increasing in $j$: for $j \le 3$ we have $2^{j} \le 8 < 9$,
while for $j \ge 4$ we have $2^{j} \ge 16 > 9$.  Hence $2^{j} \ne 9$ for every
$j \ge 0$, so $-9 = -2^{j}$ is impossible as well; equivalently
$\Delta(3,3) = -9 \notin \{\pm 2^{j}\}$.  Clause (C2) therefore fails.

The failure is not an isolated coincidence.  Iterating over the full box
$x, y \in [3, 499]$ (that is $497^{2} = 247\,009$ pairs) in
\texttt{reproduce.py}, check~[3], we find:
\[
  \#\{\, (x, y) : \Delta(x, y) = \pm 2^{j} \text{ for some } j \,\} = 39,
  \qquad
  \#\{\, (x, y) : \Delta(x, y) \ne \pm 2^{j} \text{ for all } j \,\}
  = 246\,970 .
\]
So $246\,970$ of $247\,009$ pairs -- more than $99.98\%$ -- violate the
clause.  Besides $(3,3)$ with $\Delta = -9$, small explicit violators include
$(3,4)$ with $3 \xor 4 = 7$, $3 \cdot 4 = 12$, $\Delta = -5$, and $(3,6)$ with
$3 \xor 6 = 5$, $3 \cdot 6 = 18$, $\Delta = -13$; neither $-5$ nor $-13$ is
$\pm 2^{j}$.

\begin{remark}
Since by Theorem~\ref{thm:none} there are no solutions with $x, y \ge 1$, the
set of ``other candidate pairs'' in (C2) is the entire box $x, y > 2$; the
clause is not saved by excluding solutions, because there are none to exclude.
\end{remark}

\section{Formal verification in Lean~4}\label{sec:lean}

The three refutations are formalised in core Lean~4 (toolchain
\texttt{leanprover/lean4:v4.33.1}, \texttt{import Std}, no Mathlib, no
\texttt{sorry}, no custom \texttt{axiom}, no \texttt{native\_decide}); see
\texttt{lean4/Main.lean} and the axiom audit \texttt{lean4/Check.lean}.

\begin{itemize}
  \item \texttt{two\_nim\_two} : $(2 : \mathbb{N}) \xor 2 = 0 \wedge
    2 \cdot 2 = 4$, and \texttt{not\_a\_solution} :
    $\neg\,(2 \xor 2 = 2 \cdot 2)$, both by \texttt{decide}.  This is
    Section~\ref{sec:twotwo}.
  \item \texttt{solutionsUpTo}, \texttt{only\_zero\_solution} : the bounded
    search over $0 \le x, y < 200$ equals $[(0,0)]$, and
    \texttt{no\_pos\_solutions} : the search over $1 \le x, y < 200$ is empty.
    This is the computational shadow of Theorem~\ref{thm:none}
    (Section~\ref{sec:none}).
  \item \texttt{nine\_not\_pow}, \texttt{nine\_not\_pow\_neg} :
    $\neg \exists j,\ (-9 = 2^{j})$ and $\neg \exists j,\ (-9 = -2^{j})$, and
    \texttt{clause2\_fails\_at\_3\_3}, which records $3 \xor 3 = 0$,
    $3 \cdot 3 = 9$ and both non-representability facts.  This is
    Section~\ref{sec:clause2}.
  \item \texttt{conjecture\_00000001283\_false} collects all of the above.
\end{itemize}

The axiom audit reports that every theorem depends only on the core axioms
\texttt{propext} and \texttt{Quot.sound} (never \texttt{sorryAx} or
\texttt{Lean.ofReduceBool}).

\section{Conclusion}

Conjecture 00000001283 is \textbf{false}.  The pair $(2,2)$ it lists as a
solution is not one ($2 \xor 2 = 0 \ne 4$); there is in fact no solution with
$x, y \ge 1$ whatsoever (Theorem~\ref{thm:none}); and the second clause
fails already at $(3,3)$, where the difference is $-9$, with $246\,970$ of the
$247\,009$ pairs in $[3,499]^{2}$ violating it.  These are three independent
defects, any one of which suffices to refute the statement as written.

\end{document}
